\scripture{Cap. 9, 11 – 22.}

\texteparacol{\lettrine{D}{e} Epístola beáti Pauli Apóstoli ad Hebrǽos.

\lettrine{F}{e}stinémus íngredi in illam réquiem: ut ne in idí\-psum quis íncidat incredulitátis exémplum.
Vivus est enim sermo Dei, et éfficax et penetrabílior o\-mni gládio ancípiti: et pertíngens usque ad divisiónem ánimæ ac spíritus, compágum quoque ac medullárum, et discrétor cogitatiónum et intentiónum cordis.
Et non est ulla creatúra invisíbilis in conspéctu ejus: ó\-mnia autem nuda et apérta sunt óculis ejus, ad quem nobis sermo.
Habéntes ergo Pontíficem magnum, qui penetrávit cælos, Jesum Fílium Dei: teneámus confessiónem.
Non enim habémus Pontíficem, qui non possit cómpati infirmitátibus nostris: tentátum autem per ómnia pro similitúdine absque peccáto.}{english}{\noindent Let us hasten therefore to enter into that rest; lest any man fall into the same example of unbelief. For the word of God is living and effectual, and more piercing than any two-edged sword; and reaching unto the division of the soul and the spirit, of the joints also and the marrow, and is a discerner of the thoughts and intents of the heart. Neither is there any creature invisible in his sight: but all things are naked and open to his eyes, to whom our speech is. Having therefore a great high priest that hath passed into the heavens, Jesus the Son of God: let us hold fast our confession. For we have not a high priest, who can not have compassion on our infirmities: but one tempted in all things like as we are, without sin.}